% Generated draft: requires mathematical and source review.
\documentclass[11pt]{article}
\usepackage[margin=25mm]{geometry}
\usepackage{fontspec}
\setmainfont{Latin Modern Roman}
\usepackage{amsmath,amssymb,mathtools,unicode-math}
\setmathfont{Latin Modern Math}
\usepackage{enumitem,needspace,xurl,hyperref}
\setlength{\parindent}{0pt}
\setlength{\parskip}{5pt}
\setlength{\emergencystretch}{3em}
\newcommand{\N}{\mathbb{N}}
\newcommand{\Z}{\mathbb{Z}}
\newcommand{\Q}{\mathbb{Q}}
\newcommand{\R}{\mathbb{R}}
\newcommand{\C}{\mathbb{C}}
\newcommand{\F}{\mathbb{F}}
\newcommand{\A}{\mathbb{A}}
\newcommand{\PP}{\mathbb{P}}
\newcommand{\E}{\mathbb{E}}
\newcommand{\eps}{\varepsilon}
\newcommand{\dd}{\,\mathrm{d}}
\begin{document}
\section*{UnsolvedMath ID: 30001889}
\textbf{Source problem code:} catalogue entry
\subsection{Definitions and mathematical statement}
The supplied record does not define the function \(r\), the domain of its argument \(m\), or the terms and indexing convention of OEIS A059522. Consequently, it does not determine a fully specified mathematical proposition.

The recorded question asks whether OEIS A059522 is a termwise lower bound for \(r(m)\) for every admissible \(m\). Schematically, if \(I\) denotes the intended domain and \(a_m\) denotes the term of A059522 intended to correspond to \(m\), the question is whether
\[
\forall m\in I,\qquad a_m\le r(m).
\]
Here \(I\), \(r\), and the correspondence defining \(a_m\) remain unspecified; this display is a template for the question, not a completed formulation. In particular, no starting index or boundary cases can be inferred from the supplied text.

The supplied background describes an unresolved identification with A059522, whereas the explicit statement asks only for a lower bound. Equality and a lower bound are different targets; their intended relationship requires clarification.
\subsection{Short English statement}
Does each appropriately indexed term of OEIS A059522 give a lower bound for \(r(m)\), for every admissible \(m\)? The supplied record leaves \(r\), its domain, and the sequence indexing undefined, so a precise formulation requires additional source information.
\subsection{Sources}
\begin{itemize}
\item Local UnsolvedMath registry, permanent ID \(30001889\), source problem code \texttt{catalogue entry}, title ``Sequence Bounds for the Function r of m.'' No proposer is named; the recorded proposed year is 2011.
\item Supplied report citation: \emph{Discrete Geometry (2012)}, Oberwolfach Reports, DOI 10.4171/owr/2011/44, pp.~2459--2548. \url{https://doi.org/10.4171/owr/2011/44}. The background describes Oberwolfach Reports as documenting mathematical workshops and their open-problem discussions.
\item Supplied literature source labeled ``OWR 44 (2012)'': \url{https://doi.org/10.4171/owr/2011/44}.
\item Supplied literature source: OEIS A059522, \url{https://oeis.org/A059522}. Its terms and definition are not reproduced in the record.
\item Literature assessment dated exactly 2026-08-21: supplied dated curation, not independently verified. It labels the problem open and states that the report gives only first terms and asymptotic bounds for \(r(m)\), while leaving the A059522 identification unresolved. The supplied curation lists the report and OEIS entry as checked sources and expressly requests independent verification of time-sensitive status and cited claims.
\item The registry labels the statement \texttt{corrected\_verified}; its accompanying verification describes only placing the function and its argument in mathematical delimiters without changing the question.
\end{itemize}
\end{document}
